%======================================================================
\section{Information: calibrating one expiry}
\label{sec:calibration}

Eight sections built a space in which every point is an arbitrage-free
slice.  Choosing the right point is now a plain nonlinear least-squares
problem --- which was the design goal --- and the interesting questions
are informational: what do the quotes determine, what does the objective
choose, and how are the two kept from masquerading as one another.

\subsection{The objective}

Given quotes $\sigma_i$ at log-moneyness $k_i$ for one expiry, with total
variances $w_i=\sigma_i^{2}\tau$, the calibrator solves
\begin{equation}
\min_{\theta}\;
\sum_i \left(
\frac{c(k_i;\theta)-\Black(k_i,w_i)}
     {\partial_\sigma\Black(k_i,w_i)+\eta}
\right)^{2}
\;+\;\alpha\sum_{n\ge4}n^{2\nu}a_n^{2},
\label{eq:objective}
\end{equation}
an unconstrained problem over the chart coordinates of
\cref{sec:computation}, solved by a trust-region least-squares method
\citep{BranchColemanLi1999,SciPy2020}.  Three design decisions are baked
into \eqref{eq:objective}, each with a reason.

\emph{Prices, not volatilities, in the residual.}  The model's trial
point is always a genuine law, so its prices always exist; inverting
them to implied volatility inside the loop would add a root-solve per
quote per iterate and would fail exactly on the wild trial points a
trust-region method is entitled to probe.  Fitting in price space keeps
every iterate well-defined.

\emph{Vega normalization.}  Dividing the price error by the Black vega
$\partial_\sigma\Black=\phin(d_+)\sqrt{\tau}$ at the quote restores, to
first order, a volatility error --- so the objective's units are the units
the desk thinks in, and residuals of different strikes are commensurable.
The floor $\eta$ (a small constant, \cref{tab:defaults}) matters in the
far wings, where vega vanishes: below the floor, the residual degrades
gracefully from a vol error to a scaled price error.  This is a deliberate
trade --- a vol metric at zero vega amplifies noise without bound --- but
it means short-dated wing quotes are fitted under a price-error regime,
and a reader of such fits should know it.

\emph{An explicit ridge, sparing the workhorses.}  The term
$\alpha\sum_{n\ge4}n^{2\nu}a_n^{2}$ (defaults in \cref{tab:defaults};
$\nu=1$) suppresses high-order oscillation where quotes are sparse.  The
first two body modes $a_2$, $a_3$ --- smile convexity and asymmetry, which
liquid quotes always pin --- are deliberately unpenalized.  The ridge is a
\emph{convention}, disclosed as such: it is the tie-breaker wherever the
likelihood is flat, and \cref{sec:tails} already identified the place the
tie-breaking shows up (the tail parameters).

One more residual row rides along in every fit: a smooth barrier
$\log\big(1+e^{50(\lambda_+-0.90)}\big)$, negligible below
$\lambda_+\approx0.9$ and rising steeply toward the wall.  In the
production chart the wall is unreachable
(\cref{sec:computation}), so the row is not a feasibility device; it is an
economic prior against fits that quote a nearly infinite forward-moment
budget, and it is listed here because it is part of the objective, not of
the geometry.  Its units are frankly conventional: the row is
dimensionless with unit weight, and its constants are sized so that it
is invisible below the center and reaches order one --- comparable to a
multi-vol-point data residual --- within a few hundredths of $\lambda_+$
beyond it; center and steepness are listed with the other disclosed
conventions in \cref{tab:defaults}.  A related smoothness note belongs
here: the hinge residuals of \eqref{eq:band} have a kink at the band
edge, but the solver minimizes the sum of \emph{squares}, and
$x\mapsto\big((x)^{+}\big)^{2}$ is continuously differentiable, so the
objective the trust region sees is $C^{1}$ at the edge --- the stronger
$\cos^{2}$ smoothing applied to the calendar rows of \cref{sec:calendar}
serves their strike-dependent \emph{weights}, which would otherwise
imprint kinks of their own.

\subsection{Fit targets: mid, band, haircut}

A quote is not a number; it is a bid and an ask.  The calibrator supports
three targets.

\emph{Mid.}  The residual of \eqref{eq:objective} against the mid
volatility.  Right when spreads are tight and symmetric.

\emph{Band.}  Replace the residual by a hinge on the band edges plus a
weak mid anchor: writing $[p_i^{\mathrm{lo}},p_i^{\mathrm{hi}}]$ for the
quote's price band (the vol band mapped through Black once, before the
solve) and $\phi_i$ for the vega normalizer of \eqref{eq:objective}, the
quote contributes
\begin{equation}
\frac{\big(c(k_i;\theta)-p_i^{\mathrm{hi}}\big)^{+}
+\big(p_i^{\mathrm{lo}}-c(k_i;\theta)\big)^{+}}{\phi_i}
\qquad\text{and}\qquad
\sqrt{0.05}\;\frac{c(k_i;\theta)-p_i^{\mathrm{mid}}}{\phi_i} :
\label{eq:band}
\end{equation}
zero cost anywhere inside the band, quadratic cost outside, plus five
percent of the mid-fit pull to break ties among interior solutions.  Band
fitting says honestly that the market quoted an interval.

\emph{Haircut} (the default for the examples of \cref{sec:examples}).
Bid--ask bands on listed options are conservative: the touch is often
wider than where business is done.  The haircut target shrinks each band
by a fixed volatility margin $h$ before applying \eqref{eq:band}:
\begin{equation}
\sigma_i^{\mathrm{lo}}=\min\big(\sigma_i^{\mathrm{bid}}+h,\ \sigma_i^{\mathrm{mid}}\big),
\qquad
\sigma_i^{\mathrm{hi}}=\max\big(\sigma_i^{\mathrm{mid}},\ \sigma_i^{\mathrm{ask}}-h\big),
\label{eq:haircut}
\end{equation}
with $h$ half a volatility point by default.  The $\min$/$\max$ guards
supply the degenerate-band behavior: a quote whose spread is already
tighter than $2h$ collapses to $\sigma^{\mathrm{lo}}=\sigma^{\mathrm{hi}}
=\sigma^{\mathrm{mid}}$ --- a mid fit on that strike --- so the haircut
can only tighten a band, never invert it.  The degeneracy clause is not
hypothetical; it is the live state of half the paper's data.  In the
frozen snapshot, \emph{none} of the deep-dive SPY node's bands survives
the shrink --- the live-band share is \MacSpyDecBandLivePct\% ---
which is exactly the statement that every one of its quoted spreads is
tighter than $2h=100$ vol bp (for color: the median spread is
\MacSpyDecMedianSpreadBp\ vol bp, an order of magnitude inside the
guard).  The entire SPY book fits, in effect, to mid, and the SPY
figures of \cref{sec:examples} accordingly show raw bid--ask quotes
rather than haircut bands.  The NVDA book is the other regime --- median
spreads near $2h$ (\MacNvdaLongMedianSpreadBp\ vol bp on the long node),
with \MacNvdaOneDayBandLivePct\% and \MacNvdaLongBandLivePct\% of bands
live on the featured nodes --- so NVDA is where the band mechanics
actually engage.  One recipe, two regimes, resolved per quote by
\eqref{eq:haircut} with no user decision: that is the intended design.

\subsection{What is pinned, and what is chosen}

The informational accounting of a fit, stated once, plainly.  A liquid
strip pins: the body of the distribution between the outermost quotes ---
equivalently the call curve there --- the three ATM handles, the
var-swap-style low moments, and the wing behavior \emph{through} the
quoted range.  The objective chooses: everything beyond the outermost
quotes (via the endpoint speeds, i.e.\ via the basis and the two tail
coordinates), the partition of any slack inside bid--ask bands (via the
mid anchor), and the smoothness of the body at scales below the quote
spacing (via the ridge and the order).  The choices are few, explicit,
and orthogonal to validity --- moving any of them yields a different
\emph{admissible} slice.  The tail-parameter honesty of \cref{sec:tails}
is this accounting applied to $\lambda_\pm$; the order-control study of
\cref{sec:examples} (\cref{fig:order_control}) is the same accounting
applied to $N$: at matched quote error, the extra order buys density
resolution, and the choice of $N$ decides how much of the inter-quote
structure the fitter is willing to assert.

\subsection{The order guard}
\label{sec:orderguard}

Order is a convention with a failure mode, so it is governed.  The family
at order $N$ has $N+1$ parameters; a strip of $n_q$ quotes fits at
\emph{effective} order
\begin{equation}
N_{\mathrm{eff}}
=\max\Big(\min\big(N,\ \lfloor n_q/2\rfloor-1\big),\ \min(N,6)\Big),
\label{eq:orderguard}
\end{equation}
that is: parameters are capped at half the quote count, with a floor at
order $6$ (or $N$ itself if smaller) so thin strips still fit a usable
smile.  The principled motivation is \emph{identification}: as the
parameter count approaches the quote count, the fit crosses from
estimating a curve to interpolating the noise, and the delta-method
uncertainty bands on the fitted smile blow up accordingly --- the guard
keeps the fit on the estimating side.  A second, latency-side
motivation is real but \emph{book-dependent}, and honesty requires
reporting both halves of the evidence.  On the snapshot's own thinnest
strip --- the \MacGuardQuoteCountZeroDte-quote shortest NVDA node,
\MacNvdaOneDayDays\ calendar days from the reference date
--- deliberately over-parameterizing past the guard produces no
distress at all: sweeping the parameter-to-quote ratio through
\MacCliffLoRatio\ / \MacCliffMidRatio\ / \MacCliffHiRatio\ costs
\MacCliffLoEvals\ / \MacCliffMidEvals\ / \MacCliffHiEvals\ residual
evaluations (\MacCliffLoMs, \MacCliffMidMs, \MacCliffHiMs): flat.  On
stressed same-day books, however, the reference implementation has
historically observed the same sweep multiply evaluation counts by two
orders of magnitude and push a fit from tens of milliseconds to
seconds --- an error-bar-saturated strip hands the trust region a
nearly flat valley to wander, and the failure appears without warning
on some books and not others.  The guard \eqref{eq:orderguard} is
therefore cheap insurance whose measured cost on a clean book is nil:
it holds the NVDA strip at
$N_{\mathrm{eff}}=\MacGuardEffOrderZeroDte$, and
\cref{sec:examples} shows the guarded fit on live data.

\subsection{Starting points}

Cold starts use the exactly solvable slice of \cref{sec:constspeed}:
$a_n=0$, $L=R=\log s$ with $s=\sqrt{3w_0}/\pi$ matched to the
ATM-interpolated quote variance --- a seed whose bias is known
($\approx8\%$ low in ATM variance, \cref{fig:exact}) and gone after one
iteration.  Surface sweeps warm-start each expiry from its
already-fitted shorter neighbor, since adjacent maturities carry nearly
the same smile.  A warm start is a prior in disguise --- around a
scheduled event, yesterday's unimodal seed can shepherd a genuinely
bimodal fit into the wrong basin --- so the practice is a convention to
be disclosed, and multi-start protocols are the audit for it.  The audit
was run for this paper: \MacMultiStartCount\ randomized starts on each
of \MacMultiStartNodes\ live nodes (the deep-dive SPY node and the long
NVDA node of \cref{sec:examples}) found \MacMultiStartBasins\ basin per
node --- worst intra-cluster parameter spread
\MacMultiStartWorstDTheta, worst spread in fit quality
\MacMultiStartWorstDIvBp\ vol bp.  A measurement on two liquid strips,
not a theorem, and stressed or event-laden boards remain exactly where
the caution applies.
